\documentclass[popl.tex]{subfiles}
\begin{document}

\section{Symbols at a glance}

Below we provide an inexhaustive listing of some common notation used throughout this paper.

\begin{table}[!h]
\centering
\begin{tabular}{c|l}
\hline
Notation & Meaning \\ \hline
$G=\langle \Sigma, V, P,S \rangle$ & CFG with terminals, $\Sigma$, nonterminals, $V$, productions, $P$, and start symbol $S$. \\
$A=\langle Q,\Sigma, \delta, q_\alpha, F \rangle$ & Automaton with states, $Q$, transitions, $\delta$, start state, $q_\alpha$ and final states, $F$. \\
$\err\sigma, \sigma^\star$ & Syntactically invalid string with known target language, and ground-truth repair. \\
$|\sigma|$ & Length (number of terminals) of string, $\sigma$. \\
$G'$ & Chomsky Normal Form (CNF) grammar. \\
$G_\cap$ & Intersection grammar formed by intersecting an automaton with a CFG. \\
$\ell_\cap, \mathcal{L}(G_\cap)$ & Intersection language generated by some $G_\cap$. \\
$L(\err\sigma, k)$ & Levenshtein automaton of radius $k$ for a broken string, $\err\sigma$. \\
${\color{orange}[\ldots]}$ & Orange text is related to the symbolic predicate in the Levenshtein automaton.\\
$d_{\max}$ & Maximum permitted Levenshtein edit distance (repair radius). \\
$M$ & Matrix encoding the product construction $\mathcal{L}(G)\cap\mathcal{L}\big(L(\err\sigma, k)\big)$. \\
$P@k$ & Precision at rank $k$ evaluation metric.\\
\end{tabular}\vspace{-2cm}
\end{table}
\ifSubfilesClassLoaded{\par\clearpage\bibliography{../bib/acmart}}{}
\end{document}
